\documentclass[11pt,a4paper]{article}
\usepackage[T1]{fontenc}\usepackage[utf8]{inputenc}\usepackage{lmodern}
\usepackage[margin=23mm,headheight=15pt]{geometry}
\usepackage{amsmath,amssymb,amsthm,mathtools,booktabs,longtable,array,microtype,xcolor}
\usepackage{hyperref,fancyhdr,listings,enumitem}
\definecolor{ink}{HTML}{163348}\definecolor{accent}{HTML}{147D83}
\hypersetup{colorlinks=true,linkcolor=ink,urlcolor=accent,citecolor=accent,pdfauthor={Sébastien Palcoux}}
\pagestyle{fancy}\fancyhf{}\fancyhead[L]{\sffamily\small FUSION RING CENSUS}
\fancyhead[R]{\sffamily\small Rank \Rank}\fancyfoot[L]{\sffamily\footnotesize Computational companion \textbullet\ September 2026}
\fancyfoot[R]{\sffamily\small\thepage}\renewcommand{\headrulewidth}{0.4pt}
\newtheorem{theorem}{Theorem}[section]\newtheorem{proposition}[theorem]{Proposition}
\newtheorem{lemma}[theorem]{Lemma}\newtheorem{corollary}[theorem]{Corollary}
\theoremstyle{definition}\newtheorem{definition}[theorem]{Definition}\newtheorem{remark}[theorem]{Remark}
\newcommand{\ZZ}{\mathbb Z}\newcommand{\CC}{\mathbb C}\newcommand{\Q}{\mathbb Q}
\newcommand{\NN}{\mathbb Z_{\geq0}}\newcommand{\FPdim}{\operatorname{FPdim}}
\setlength{\parindent}{0pt}\setlength{\parskip}{0.45em}\setlength{\emergencystretch}{3em}
\lstset{basicstyle=\ttfamily\footnotesize,breaklines=true,columns=fullflexible,frame=single,rulecolor=\color{accent},aboveskip=1em,belowskip=1em}
\setcounter{tocdepth}{2}
\newcommand{\Rank}{4}
\begin{document}
\thispagestyle{empty}{\sffamily\small\color{accent} EXACT ENUMERATION / DATA / VERIFICATION}\par\vspace{1cm}
{\sffamily\Huge\bfseries\color{ink} Rank-4 fusion rings\par}
\vspace{0.35cm}
{\sffamily\Large A reproducible census through multiplicity 128\par}
\vspace{0.65cm}{\large Sébastien Palcoux}\par{\small BIMSA}\par\vspace{0.35cm}{\small Computational companion, 23 September 2026}\par\vspace{0.8cm}
\begin{abstract}
We explain the exact enumeration of all based fusion rings of rank 4 through multiplicity 128, comprising 889,530 isomorphism classes. The account separates the mathematical coverage argument, the completed search, and independent verification of the emitted tensors. All duality types are included, and no categorifiability filter is applied. Code, machine-readable counts, full multiplication tensors and provenance records accompany the text.
\end{abstract}\vspace{0.4cm}\noindent\textbf{AI assistance and verification.} GPT-6 Astra Pro assisted with code, exposition and computational checks. Exhaustive-search arguments and exact independent tensor verification are recorded separately; no proof-assistant certification or peer-reviewed acceptance is claimed.\vfill\noindent\textbf{Scope.} A census of based rings, not a classification of tensor categories. Only completed whole-rank results are released.\newpage\tableofcontents\newpage

\section{Objects, equivalence and the counting convention}
A \emph{fusion ring} here is a free abelian group with a distinguished finite basis
$B=\{b_0,b_1,\ldots,b_{r-1}\}$, an associative multiplication with unit $b_0$,
and an involution $i\mapsto i^*$ fixing $0$. Its structure constants satisfy
\[
b_i b_j=\sum_k N_{ij}^{k}b_k,\qquad N_{ij}^{k}\in\NN,
\]
\begin{align}
N_{0j}^{k}=N_{j0}^{k}&=\delta_{jk},&N_{ij}^{0}&=\delta_{i^*,j},\\
N_{ij}^{k}&=N_{i^*k}^{j}=N_{kj^*}^{i}.&&\label{eq:reciprocity}
\end{align}
These conventions agree with the usual based-ring definition; see
\cite{EGNO,VS}. No assumption of categorifiability, commutativity or integrality
of Frobenius--Perron dimensions is imposed unless proved in the relevant rank.

A based isomorphism is a bijection of the distinguished bases preserving the
unit and multiplication. In tensor notation it is a permutation $q$ fixing $0$
such that
\[
 N_{ij}^{k}={N'}_{q(i)q(j)}^{q(k)}\quad\hbox{for every }i,j,k.
\]
The coefficient of the unit determines the involution, so such an isomorphism
automatically intertwines duality. The \emph{multiplicity} is the maximum entry
of the \emph{full} ordered multiplication tensor. It is at least one, including
for group rings. Put
\[
 c_r(m)=\#\{\text{rank-}r\text{ fusion rings of multiplicity exactly }m\}/\cong,
 \qquad C_r(M)=\sum_{m=1}^{M}c_r(m).
\]
A list through $M$ contains all rings of multiplicity at most $M$, not only
those of multiplicity equal to $M$. This distinction is essential when exporting
OEIS b-files.

\subsection{Why duality types can be searched separately}
An involution on the $r-1$ nonunit elements consists of $p$ disjoint transpositions
and $r-1-2p$ fixed points, with $0\leq p\leq\lfloor(r-1)/2\rfloor$. Fix one
representative $d_p$ of each type. An isomorphism between two tensors with
involution $d_p$ commutes with $d_p$. Hence it is sufficient to quotient that
stratum by the centralizer
\[
 Z_{S_{r-1}}(d_p),\qquad |Z_{S_{r-1}}(d_p)|=2^p p!(r-1-2p)!.
\]
Different $p$ cannot be isomorphic. Therefore a whole-rank census is complete
exactly when every required stratum has been searched exhaustively. Finishing
only the non-self-dual strata does not certify a whole-rank counting term.
\section{The certified release at rank 4}
The released bound is $M=128$ and the cumulative number of classes is $C_{4}(128)=889,530$. Every duality type is included. The accompanying full tensor file has one record per based-isomorphism class.


\begin{center}\begin{tabular}{rr}\toprule Exact multiplicity $m$ & $c_{\Rank}(m)$\\\midrule
1 & 10\\
2 & 17\\
3 & 24\\
4 & 45\\
5 & 55\\
6 & 81\\
7 & 92\\
8 & 137\\
9 & 151\\
10 & 186\\
11 & 238\\
12 & 291\\
13 & 246\\
14 & 340\\
15 & 349\\
16 & 525\\
\bottomrule\end{tabular}\end{center}
The complete table through the released bound appears in Appendix A.

The data and their checksums are recorded in \path{results/census.json}. The corresponding completed-run evidence and independent audit are under \path{verification/rank4/}.

\section{Why the rank-four specialization is commutative}
The coefficient-of-unit functional $\tau$ gives the positive inner product
$\langle x,y\rangle=\tau(xy^*)$ on the complexified based ring. The basis is
orthonormal and multiplication is closed under adjoints. Consequently its
complexification is semisimple. The Frobenius--Perron dimension character is a
one-dimensional representation. A noncommutative semisimple complex algebra
with a one-dimensional representation has dimension at least $1+2^2=5$.
Thus a rank-four fusion ring is commutative. This justifies the specialization
below without any categorical assumption.
\section{An independent rank-four census}\label{sec:census}
We give a complete integer enumeration, with no input from an external
catalogue. In this section all parameters lie in $\{0,\ldots,M\}$.
The unit coefficients ensure that the multiplicity is the maximum of one
and the parameters. Commutativity was proved above; duality has just two possible types.

\subsection{Self-dual parameters and elimination}
Frobenius reciprocity gives the following general matrices:
\[
N_x=\begin{pmatrix}0&1&0&0\\1&a&b&c\\0&b&d&e\\0&c&e&f\end{pmatrix},\quad
N_y=\begin{pmatrix}0&0&1&0\\0&b&d&e\\1&d&g&h\\0&e&h&i\end{pmatrix},\quad
N_z=\begin{pmatrix}0&0&0&1\\0&c&e&f\\0&e&h&i\\1&f&i&j\end{pmatrix}.
\]
Their pairwise commutativity is equivalent to the six equations
\begin{align}
 ad-b^2+bg+ch-d^2-e^2+1&=0,\label{eq:census1}\\
 ae-bc+bh+ci-de-ef&=0,\label{eq:census2}\\
 be-cd+dh-eg+ei-fh&=0,\label{eq:census3}\\
 af+bi-c^2+cj-e^2-f^2+1&=0,\label{eq:census4}\\
 bf-ce+di-eh+ej-fi&=0,\label{eq:census5}\\
 df-e^2+gi-h^2+hj-i^2+1&=0.\label{eq:census6}
\end{align}
These are necessary and sufficient for associativity. Indeed the first
row of $N_uN_v$ records $uv$. Commutativity implies that every other row
of $N_uN_v-\sum_wN_{uv}^wN_w$ is obtained by multiplying its zero first
row by a regular matrix, so the difference vanishes.

The six mixed coordinates $b,c,d,f,h,i$ form one orbit under
permutations of $x,y,z$. Every based-ring orbit therefore has a
representative with
\[
 b=\min\{b,c,d,f,h,i\}.
\]
For $b>0$, enumerate $b,c,d,e,f$, with $c,d,f\geq b$, and put
\begin{align*}
 H_0&=e(d+f)+bc,&G_0&=b^2+d^2+e^2-1,\\
 J_0&=c^2+e^2+f^2-1,&
 A&=e(b^2-c^2)+bc(f-d),\\
 C&=e(bf-ce),&
 B_0&=b^3e-b^2cd+b(d-f)H_0-ebG_0+ecH_0.
\end{align*}
The first, second and fourth equations give
\[
 h=\frac{H_0-ae-ci}{b},\qquad
 g=\frac{G_0-ad-ch}{b},\qquad
 j=\frac{J_0-af-bi}{c}.
\]
After substitution, the third equation multiplied by $b^2$ is
\begin{equation}\label{eq:census-linear}
 Ca+Ai+B_0=0.
\end{equation}
If $A\ne0$, let $q=\gcd(A,C)$. There is no solution unless $q\mid B_0$.
Otherwise $a$ runs through the unique residue class solving
$(C/q)a\equiv-B_0/q\pmod{|A/q|}$, and $i$ is determined.
A modulus of one means every bounded $a$. If $A=0\ne C$, determine
$a$ and enumerate $i$. If $A=C=0$, require $B_0=0$ and retain all bounded
$a,i$. Impose $0\leq a,g,j\leq M$ and $b\leq h,i\leq M$,
exact divisibility, and all six original equations.

If $b=0<c$, determine
\[
 h=\frac{d^2+e^2-ad-1}{c},\quad
 i=\frac{e(d+f-a)}c,\quad
 j=\frac{c^2+e^2+f^2-af-1}c.
\]
For $e>0$, integrality of $i$ gives
$a\equiv d+f\pmod{c/\gcd(c,e)}$. Enumerate $c,d,e,f$ and this
progression for $a$; the bounds on $h,i,j$ give further linear interval
restrictions. The equation $eg=h(d-f)+ei-cd$ determines $g$.
For $e=0$, nonnegativity forces $d,h>0$ and
\[
 \gcd(c,d)=1,\qquad a\equiv d-d^{-1}\pmod c,\qquad
 f=d-\frac{cd}{h}.
\]
Indeed $ch=d^2-ad-1$ implies the coprimality, and $h(d-f)=cd$
gives the last identity. Enumerate $c,d,a$, determine $h,f,j$, and
retain every bounded $g$. The congruence modulo one is unrestricted.

Finally, $b=c=0$ gives
\[
 ad-d^2-e^2+1=0,\quad e(a-d-f)=0,\quad af-e^2-f^2+1=0.
\]
For $e>0$ these are $a=d+f$ and $df=e^2-1$. Enumerate the bounded
factor pairs, including $e=1,d=0$ with every bounded $f$, then enumerate
$h,i$ and determine
\[
 eg=h(d-f)+ei,\qquad ej=eh+i(f-d).
\]
For $e=0$, the preliminary equations force $a=0,d=f=1$. The remaining
equation is $gi+hj=h^2+i^2-2$. If $h>0$, enumerate $h,i,g$ and solve
for $j$; if $h=0<i$, solve for $g$ and retain every bounded $j$.
The case $h=i=0$ is impossible. In every branch, all six equations and
all coefficient bounds are tested again.

\subsection{One dual pair}
If $X^*=X$ and $Y^*=Z$, the general matrices are
\[
N_X=\begin{pmatrix}0&1&0&0\\1&a&b&b\\0&b&c&d\\0&b&d&c\end{pmatrix},\quad
N_Y=\begin{pmatrix}0&0&1&0\\0&b&c&d\\0&d&e&f\\1&c&e&e\end{pmatrix},\quad
N_Z=N_Y^{\mathsf T}.
\]
Associativity is equivalent to
\begin{align*}
 ac-b^2+2be-c^2-d^2+1&=0,&ad-b^2+be+bf-2cd&=0,\\
 b(c-d)+d(e-f)&=0,&-c^2+d^2-e^2+f^2-1&=0.
\end{align*}
Enumerate $c,d,e$. The last equation determines the unique nonnegative
candidate $f=\sqrt{1+c^2-d^2+e^2}$; retain it only if integral and
bounded. If $c\ne d$, determine $b=-d(e-f)/(c-d)$. If $c=d$, require
$d(e-f)=0$ and retain every bounded $b$. Determine $a$ from the first
equation if $c>0$, or from the second if $d>0$; if $c=d=0$, retain
all bounded $a$. Finally test all four equations and all bounds.
The square-root implementation corrects its initial estimate by integer
comparisons and checks its square exactly.

\begin{proposition}[Completeness and cost]\label{prop:census}
The preceding algorithm enumerates all rank-four fusion rings of
multiplicity at most $M$, up to based-ring isomorphism. The parameter
search uses $O((M+1)^6)$ integer operations before canonicalization.
\end{proposition}
\begin{proof}
The two presentations exhaust duality and reciprocity. Each branch either
solves a necessary linear equation by exact division or exhausts every
allowed value when its coefficient vanishes. Final equation checks are
necessary and sufficient for associativity. Thus no ring is omitted or
spuriously admitted. In the self-dual search choose the lexicographically
least parameter tuple among its permutations satisfying the same
minimal-$b$ condition. Comparing only those images is essential: the
unconstrained least tuple need not satisfy that condition. Each orbit has
exactly one retained tuple. In the other duality type the only possible
nontrivial based automorphism interchanges $Y,Z$, and fixes all six
parameters. Thus every retained non-self-dual tuple is already one class.
The full-tensor interface additionally canonicalizes under all six
permutations and checks distinctness independently.

There are five outer parameters in the positive-$b$ branch. Solving
\eqref{eq:census-linear} takes $O(\log(M+1))$ arithmetic operations for
the greatest common divisor and at most $M+1$ candidates. The apparently
singular branch $A=C=B_0=0$ cannot occur there: if $e=0$, then
$A=0$ forces $f=d$ and $B_0=-b^2cd\ne0$; if $e>0$, the equations
$A=C=0$ give $cd=be$, $bf=ce$, and substitution gives $B_0=be\ne0$.
The zero-$b$ branches use at most five free coordinates, and the
non-self-dual search at most five. The stated $O((M+1)^6)$ arithmetic
bound follows. Bit costs grow with $\log(M+1)$; the implementation
restricts its integer range explicitly.
\end{proof}

\subsection{Current implementation and complete run}
The current implementation is \texttt{code/rank4\_census.cpp}; the
reproduction entry point is \texttt{code/run\_census.py}. The full
run through multiplicity $128$ produces $889530$ based-isomorphism
classes. Every output tensor is independently checked for the full
fusion-ring axioms and canonicalized over all six unit-fixing
permutations. The exact counts and execution reports are in the
accompanying release. Exhaustiveness rests on the argument above,
not on comparison with an existing list.

\section{What the independent verification establishes}
The enumerator and verifier have separate mathematical tasks. Exhaustiveness
uses the complete coordinate parameterization, the safe-pruning proof and the
normal completion of every required branch. The verifier reads only the emitted
\emph{full tensors}; it neither reconstructs the parameterization nor trusts the
generator's polynomial list.

For every tensor it checks nonnegative integral coefficients, both unit laws,
uniqueness and involutivity of the dual, both reciprocity identities, and every
integer associativity identity for all four indices. It then canonicalizes
under \emph{every} unit-fixing basis permutation. A matching hash is only a
lookup aid: actual canonical tensors are compared. This detects both literal
repetitions and different labellings of an isomorphic based ring.

Passing these checks proves soundness and uniqueness of the emitted data; by
itself it would not prove that no ring was omitted. That is why normal completion
logs and the mathematical coverage argument are retained separately. Comparison
with the original ancillary file is an additional containment/regression test,
not an input to enumeration and not a replacement for an exhaustiveness proof.

\subsection{Provenance and reproducibility}
The retained Vercleyen--Slingerland v4 ancillary file has SHA-256
\begin{center}\footnotesize\ttfamily
b587c3f683157369da7cc090f762bbff0e18dcf11acc3f5031d6bc0c643843f8
\end{center}
Its 28,451 records contain 25,138 distinct classes. The 3,313 redundant records
occur at rank five and exact multiplicities 9--12. Their original and corrected
counts are respectively
\[
(1463,1794,2283,3049)\quad\longrightarrow\quad(863,1082,1383,1948).
\]
Every deletion has an explicit basis-permutation witness in the audit folder.
This audit does not establish completeness of the source's partial-search cells.

The primary manifest \path{results/census.json} records every released count,
bound, filename and checksum. Compressed tensors use the original line-per-ring,
nested-curly-brace format. An entry is $N[i][j][k]=N_{ij}^{k}$: the output index
has already been dualized where the internal parameter convention requires it.
The raw tensor needs no additional index conversion.

\subsection{Certified datasets and reproducibility}
The primary manifest gives the complete bound, exact counts and data checksums
for every rank. Each \path{verification/rankR/} directory records the independent
audit and source identity for that dataset. For the laptop censuses it also
retains the completed enumeration, commands, environment, phase timings and
nonempty subprocess logs. Every required duality stratum completed normally;
independent canonicalization found zero duplicate classes.

The independent checks of ranks five through eight ran on GitHub using the
supplied tensors. They did not rerun enumeration. The per-rank reports identify
the exact verification revision and run, and certify agreement with the established
exact-multiplicity count prefix. See \path{verification/README.md} for the evidence
index and \path{docs/PROVENANCE.md} for attribution.

The local launcher has no clock limit, while the manual hosted workflow uses a
finite shared deadline covering every computational phase. A timeout never
certifies the unfinished bound. Only a complete, independently verified whole-rank
candidate can enter the release; a later failed attempt cannot replace it.
Reproduction commands and supported limits are in \path{docs/REPRODUCIBILITY.md}.

\section{Reading and reusing the data}
The result files count based rings, not monoidal categories or realizations of
a ring. One tensor can admit several inequivalent categorifications, or none.
No number here is a classification of fusion categories. The rank-specific
count files use exact multiplicity; cumulative totals belong only in the
summary column. Proposed OEIS additions must be regenerated from a complete
manifest, not copied from an interrupted log.

The retained code was generated with AI assistance and is accompanied by explicit
mathematical coverage arguments and independent integer checks. These checks
are not a formal proof in a proof assistant. The present text is an explanatory
computational companion, not a claim of journal acceptance or external peer review.

\par\noindent\begin{minipage}{\linewidth}
\section{Reproduction commands}
From the repository root on a local computer, request a complete run with no clock limit. An interruption does not certify a smaller census.
\par\noindent\begin{minipage}{\linewidth}
\begin{lstlisting}
python3 scripts/run_laptop.py --rank 4 --bound 128
\end{lstlisting}
\end{minipage}\par
To verify the supplied release without recomputing it:
\par\noindent\begin{minipage}{\linewidth}
\begin{lstlisting}
python3 scripts/check_release.py --full
\end{lstlisting}
\end{minipage}\par
For a shared-budget frontier search, use \path{scripts/extend_census.py}. For ranks five through eight it starts at the released bound plus one and restarts the unfinished bound without claiming checkpoint resumption. The local command above verifies all tensors and the old count prefix before packaging a result. The separate timed frontier wrapper bounds every phase. Rank three is instead capped at 1000.

\end{minipage}\par
\clearpage\appendix
\section{All exact-multiplicity rank-four counts}
\begin{longtable}{rr@{\hspace{1.1cm}}rr@{\hspace{1.1cm}}rr@{\hspace{1.1cm}}rr}
\toprule $m$ & $c_4(m)$ & $m$ & $c_4(m)$ & $m$ & $c_4(m)$ & $m$ & $c_4(m)$\\\midrule\endhead
1 & 10 & 2 & 17 & 3 & 24 & 4 & 45\\
5 & 55 & 6 & 81 & 7 & 92 & 8 & 137\\
9 & 151 & 10 & 186 & 11 & 238 & 12 & 291\\
13 & 246 & 14 & 340 & 15 & 349 & 16 & 525\\
17 & 424 & 18 & 477 & 19 & 513 & 20 & 713\\
21 & 704 & 22 & 668 & 23 & 654 & 24 & 1,121\\
25 & 868 & 26 & 938 & 27 & 989 & 28 & 1,286\\
29 & 1,356 & 30 & 1,347 & 31 & 1,220 & 32 & 1,671\\
33 & 1,460 & 34 & 1,579 & 35 & 1,613 & 36 & 2,301\\
37 & 1,646 & 38 & 1,824 & 39 & 2,111 & 40 & 2,653\\
41 & 2,131 & 42 & 2,588 & 43 & 2,169 & 44 & 2,857\\
45 & 2,699 & 46 & 2,639 & 47 & 2,452 & 48 & 3,985\\
49 & 2,947 & 50 & 3,314 & 51 & 3,378 & 52 & 3,776\\
53 & 3,250 & 54 & 3,704 & 55 & 3,785 & 56 & 5,213\\
57 & 4,182 & 58 & 3,863 & 59 & 3,987 & 60 & 6,004\\
61 & 4,383 & 62 & 4,472 & 63 & 5,150 & 64 & 6,557\\
65 & 5,107 & 66 & 5,773 & 67 & 4,791 & 68 & 6,260\\
69 & 6,001 & 70 & 6,205 & 71 & 5,799 & 72 & 8,710\\
73 & 5,858 & 74 & 6,223 & 75 & 7,280 & 76 & 9,507\\
77 & 6,953 & 78 & 7,917 & 79 & 6,933 & 80 & 10,208\\
81 & 8,031 & 82 & 7,939 & 83 & 7,408 & 84 & 10,993\\
85 & 8,464 & 86 & 8,568 & 87 & 9,088 & 88 & 11,767\\
89 & 9,145 & 90 & 10,765 & 91 & 9,896 & 92 & 11,188\\
93 & 10,385 & 94 & 10,153 & 95 & 10,048 & 96 & 15,447\\
97 & 10,003 & 98 & 11,249 & 99 & 12,082 & 100 & 14,037\\
101 & 11,361 & 102 & 13,068 & 103 & 10,861 & 104 & 16,150\\
105 & 14,958 & 106 & 12,338 & 107 & 12,102 & 108 & 16,509\\
109 & 12,849 & 110 & 14,390 & 111 & 15,119 & 112 & 18,964\\
113 & 13,293 & 114 & 16,769 & 115 & 14,900 & 116 & 17,484\\
117 & 16,252 & 118 & 14,963 & 119 & 16,058 & 120 & 25,126\\
121 & 16,103 & 122 & 15,670 & 123 & 17,292 & 124 & 19,682\\
125 & 17,760 & 126 & 20,720 & 127 & 17,015 & 128 & 23,185\\
\bottomrule\end{longtable}

\begin{thebibliography}{9}\small\setlength{\itemsep}{2pt}
\bibitem{EGNO} P. Etingof, S. Gelaki, D. Nikshych and V. Ostrik,
\emph{Tensor Categories}, Mathematical Surveys and Monographs 205,
American Mathematical Society, 2015.
\bibitem{VS} G. Vercleyen and J. Slingerland,
\emph{On Low Rank Fusion Rings}, Journal of Mathematical Physics 64 (2023),
091703; \href{https://arxiv.org/abs/2205.15637v4}{arXiv:2205.15637v4}.
\bibitem{Ostrik} V. Ostrik, \emph{On formal codegrees of fusion categories},
Mathematical Research Letters 16 (2009), 895--901;
\href{https://arxiv.org/abs/0810.3242}{arXiv:0810.3242}.
\bibitem{DP} J. Dong and S. Palcoux, rank-four and rank-five census derivations,
project manuscripts and computational supplements, September 2026.
The relevant derivations are retained in the repository's \path{methods/} directory.
\bibitem{OEIS} OEIS Foundation Inc., \emph{The On-Line Encyclopedia of Integer
Sequences}, entries A348305, A354471, A354472, A354473, A354475, A354476
and A354477. Proposed updates in this repository are not submitted edits.
\end{thebibliography}
\end{document}
